% GENERATED FILE. Edit canonical lessons, then run npm run generate:books.
% canonical-source-hash: fnv1a64:71adf2934823268e
% canonical-lessons: BN-C431-uposagor, BN-C431-khari, BN-C431-bodvip, BN-C431-poli, BN-C431-bhorat-kora

\chapter{Bays, Creeks and Deltas}
\label{ch:bn-bays-creeks-and-deltas}
\hlchaptermodality{431}

\begin{chapteropening}
\textbf{By the end of this chapter:} \emph{I can talk about where rivers meet the sea: bays, creeks, deltas and silt.}

Complete the last lesson of chapter 431: I can talk about where rivers meet the sea: bays, creeks, deltas and silt.
\end{chapteropening}

\section[\texorpdfstring{\bn{উপসাগর} (upôsāgôr)}{উপসাগর (upôsāgôr)}]{\bn{উপসাগর} (upôsāgôr) — a bay; a gulf}

\label{lesson:BN-C431-uposagor}



\begin{quote}
Before the new one: say the Bengali for to irrigate, then the Bengali for to farm.
\end{quote}

\subsection*{You'll want to know: \bn{উপসাগর}}
\textbf{\bn{উপসাগর}} — \emph{upôsāgôr} — \textquotedblleft{}a bay; a gulf\textquotedblright{}.

\begin{quote}
\bn{নদীটা} \bn{দক্ষিণে} \bn{গিয়ে} \bn{উপসাগরে} \bn{পড়ে।}
\end{quote}

\begin{quote}
\emph{nôdīṭā dôkkhine giye upôsāgôre pôṛe.}
\end{quote}

\begin{quote}
The river flows south and falls into the bay.
\end{quote}

The Bay of Bengal has its own name, built on this word.

\begin{cousinweb}
\textbf{\bn{উপসাগর}} (\emph{upôsāgôr}) joins Sanskrit \emph{upa}, \textquotedblleft{}near, lesser\textquotedblright{}, to \emph{sāgara}, \textquotedblleft{}the sea\textquotedblright{} (Monier-Williams): a lesser sea, a modern coinage for a bay.
\end{cousinweb}

\begin{grammarlens}[title={What you've built}]
One more word for getting around.
\end{grammarlens}

\subsection*{Your turn}
\begin{itemize}
\raggedright
  \item \emph{Say it:} \emph{upôsāgôr}
  \item \emph{Say it:} \emph{upôsāgôr}, once more
  \item \emph{Say it:} \emph{nôdīṭā dôkkhine giye upôsāgôre pôṛe.} — The river flows south and falls into the bay.
  \item \emph{Recall:} say the Bengali for to irrigate, then the Bengali for to farm, then say \emph{upôsāgôr} again
\end{itemize}

\begin{tcolorbox}[breakable,colback=teal!4,colframe=teal!35!black,title={Before you move on}]
What does \bn{উপসাগর} mean? (\textquotedblleft{}A bay; a gulf\textquotedblright{}.) Say it once more.
\end{tcolorbox}
\section[\texorpdfstring{\bn{খাঁড়ি} (khā̃ṛi)}{খাঁড়ি (khā̃ṛi)}]{\bn{খাঁড়ি} (khā̃ṛi) — a creek; an inlet}

\label{lesson:BN-C431-khari}



\begin{quote}
Before the new one: say the Bengali for to farm, then the Bengali for a bay.
\end{quote}

\subsection*{You'll want to know: \bn{খাঁড়ি}}
\textbf{\bn{খাঁড়ি}} — \emph{khā̃ṛi} — \textquotedblleft{}a creek; an inlet\textquotedblright{}.

\begin{quote}
\bn{জোয়ারের} \bn{সময়} \bn{খাঁড়িতে} \bn{জল} \bn{ঢোকে।}
\end{quote}

\begin{quote}
\emph{joyārer sômôy khā̃ṛite jôl ḍhoke.}
\end{quote}

\begin{quote}
At high tide the water comes into the creek.
\end{quote}

\begin{cousinweb}
\textbf{\bn{খাঁড়ি}} (\emph{khā̃ṛi}) is inherited: CDIAL lists Bengali \emph{khāṛī}, \textquotedblleft{}a creek, an inlet\textquotedblright{}, under the reconstructed \emph{khāḍa} (CDIAL 3860), and again under \emph{khaḍḍa} (CDIAL 3790).
\end{cousinweb}

\begin{grammarlens}[title={What you've built}]
One more word for getting around.
\end{grammarlens}

\subsection*{Your turn}
\begin{itemize}
\raggedright
  \item \emph{Say it:} \emph{khā̃ṛi}
  \item \emph{Say it:} \emph{khā̃ṛi}, once more
  \item \emph{Say it:} \emph{joyārer sômôy khā̃ṛite jôl ḍhoke.} — At high tide the water comes into the creek.
  \item \emph{Recall:} say the Bengali for to farm, then the Bengali for a bay, then say \emph{khā̃ṛi} again
\end{itemize}

\begin{tcolorbox}[breakable,colback=teal!4,colframe=teal!35!black,title={Before you move on}]
What does \bn{খাঁড়ি} mean? (\textquotedblleft{}A creek; an inlet\textquotedblright{}.) Say it once more.
\end{tcolorbox}
\section[\texorpdfstring{\bn{বদ্বীপ} (bôdvīp)}{বদ্বীপ (bôdvīp)}]{\bn{বদ্বীপ} (bôdvīp) — a delta}

\label{lesson:BN-C431-bodvip}



\begin{quote}
Before the new one: say the Bengali for a bay, then the Bengali for a creek.
\end{quote}

\subsection*{You'll want to know: \bn{বদ্বীপ}}
\textbf{\bn{বদ্বীপ}} — \emph{bôdvīp} — \textquotedblleft{}a delta\textquotedblright{}.

\begin{quote}
\bn{বদ্বীপে} \bn{অনেক} \bn{নদী} \bn{আর} \bn{খাল} \bn{আছে।}
\end{quote}

\begin{quote}
\emph{bôdvīpe ônek nôdī ār khāl āchhe.}
\end{quote}

\begin{quote}
There are many rivers and canals in the delta.
\end{quote}

The \bn{ব} in front is the letter: an island shaped like \bn{ব}, as the Greek delta is shaped like its letter.

\begin{cousinweb}
\textbf{\bn{বদ্বীপ}} (\emph{bôdvīp}) is usually explained as \bn{ব}, the letter, and \textbf{\bn{দ্বীপ}}, \textquotedblleft{}an island\textquotedblright{}, Sanskrit \emph{dvīpa} (Monier-Williams): a coinage that renders the Greek \emph{delta}, named for a triangular letter, with a Bengali one.
\end{cousinweb}

\begin{grammarlens}[title={What you've built}]
One more word for getting around.
\end{grammarlens}

\subsection*{Your turn}
\begin{itemize}
\raggedright
  \item \emph{Say it:} \emph{bôdvīp}
  \item \emph{Say it:} \emph{bôdvīp}, once more
  \item \emph{Say it:} \emph{bôdvīpe ônek nôdī ār khāl āchhe.} — There are many rivers and canals in the delta.
  \item \emph{Recall:} say the Bengali for a bay, then the Bengali for a creek, then say \emph{bôdvīp} again
\end{itemize}

\begin{tcolorbox}[breakable,colback=teal!4,colframe=teal!35!black,title={Before you move on}]
What does \bn{বদ্বীপ} mean? (\textquotedblleft{}A delta\textquotedblright{}.) Say it once more.
\end{tcolorbox}
\section[\texorpdfstring{\bn{পলি} (pôli)}{পলি (pôli)}]{\bn{পলি} (pôli) — silt}

\label{lesson:BN-C431-poli}



\begin{quote}
Before the new one: say the Bengali for a creek, then the Bengali for a delta.
\end{quote}

\subsection*{You'll want to know: \bn{পলি}}
\textbf{\bn{পলি}} — \emph{pôli} — \textquotedblleft{}silt\textquotedblright{}.

\begin{quote}
\bn{বন্যার} \bn{পরে} \bn{খেতে} \bn{নতুন} \bn{পলি} \bn{পড়ল।}
\end{quote}

\begin{quote}
\emph{bônyār pôre khete nôtun pôli pôṛlô.}
\end{quote}

\begin{quote}
After the flood fresh silt settled on the fields.
\end{quote}

The river mud that makes Bengal's fields rich: \bn{পলিমাটি}, silt soil.

\begin{cousinweb}
\textbf{\bn{পলি}} (\emph{pôli}) is the fine river mud left by floods; its origin is not settled.
\end{cousinweb}

\begin{grammarlens}[title={What you've built}]
One more word for getting around.
\end{grammarlens}

\subsection*{Your turn}
\begin{itemize}
\raggedright
  \item \emph{Say it:} \emph{pôli}
  \item \emph{Say it:} \emph{pôli}, once more
  \item \emph{Say it:} \emph{bônyār pôre khete nôtun pôli pôṛlô.} — After the flood fresh silt settled on the fields.
  \item \emph{Recall:} say the Bengali for a creek, then the Bengali for a delta, then say \emph{pôli} again
\end{itemize}

\begin{tcolorbox}[breakable,colback=teal!4,colframe=teal!35!black,title={Before you move on}]
What does \bn{পলি} mean? (\textquotedblleft{}Silt\textquotedblright{}.) Say it once more.
\end{tcolorbox}
\section[\texorpdfstring{\bn{ভরাট} \bn{করা} (bhôrāṭ kôrā)}{ভরাট করা (bhôrāṭ kôrā)}]{\bn{ভরাট} \bn{করা} (bhôrāṭ kôrā) — to fill in (a pond or low land)}

\label{lesson:BN-C431-bhorat-kora}



\begin{quote}
Before the new one: say the Bengali for a delta, then the Bengali for silt.
\end{quote}

\subsection*{You'll want to know: \bn{ভরাট} \bn{করা}}
\textbf{\bn{ভরাট} \bn{করা}} — \emph{bhôrāṭ kôrā} — \textquotedblleft{}to fill in (a pond or low land)\textquotedblright{}.

\begin{quote}
\bn{লোকেরা} \bn{পুরনো} \bn{পুকুরটা} \bn{ভরাট} \bn{করল।}
\end{quote}

\begin{quote}
\emph{lokerā purôno pukurṭā bhôrāṭ kôrlô.}
\end{quote}

\begin{quote}
People filled in the old pond.
\end{quote}

Past (he/she …): \bn{তিনি} \bn{জমিটা} \bn{ভরাট} \bn{করলেন}, \emph{tini jômiṭā bhôrāṭ kôrlen}.

\begin{cousinweb}
\textbf{\bn{ভরাট} \bn{করা}} (\emph{bhôrāṭ kôrā}) joins \bn{ভরাট}, \textquotedblleft{}filled up, solid\textquotedblright{}, built on \textbf{\bn{ভরা}}, \textquotedblleft{}full\textquotedblright{} (CDIAL 9397), to \bn{করা}, to do.
\end{cousinweb}

\begin{grammarlens}[title={What you've built}]
One more word for getting around.
\end{grammarlens}

\subsection*{Your turn}
\begin{itemize}
\raggedright
  \item \emph{Say it:} \emph{bhôrāṭ kôrā}
  \item \emph{Say it:} \emph{bhôrāṭ kôrā}, once more
  \item \emph{Say it:} \emph{lokerā purôno pukurṭā bhôrāṭ kôrlô.} — People filled in the old pond.
  \item \emph{Recall:} say the Bengali for a delta, then the Bengali for silt, then say \emph{bhôrāṭ kôrā} again
\end{itemize}

\begin{tcolorbox}[breakable,colback=teal!4,colframe=teal!35!black,title={Before you move on}]
What does \bn{ভরাট} \bn{করা} mean? (\textquotedblleft{}To fill in (a pond or low land)\textquotedblright{}.) Say it once more.
\end{tcolorbox}
